\scanpage{11}

\uline{\S\ 2) (Décompositions cellulaires)} \deleted{\unsure{\S\ 1}}

Soit $\mathcal P=(X_i)_{i\in I}$ une partition \inserted{finie} de $X$, et soit $F_{\mathcal P}$ l'ens des parties de $X$ qui sont des réunions \deleted{finies} de $X_i$. Alors \deleted{\unsure{le}} $F_{\mathcal P}$ satisfait les conditions a), b), et $\mathcal A_{(F_{\mathcal P})}$ est l'algèbre des fonctions sur $X$ à valeurs dans $\Lambda$ qui sont constantes sur chaque $X_i$; i.e. qui se factorisent par l'ens quotient $X/R$, où $R$ \deleted{\unsure{la}} est la relation d'équivalence associée à la partition.

\margin{$F_{\mathcal P}$ est en \unsure{corres.} biunivoque avec l'ens des parties de $X/R$.}

\margin{\uline{NB} La partition $\mathcal P$ est connue quand on connaît $F_{\mathcal P}$ \deleted{\unsure{la}} : elle est formée des \unsure{éléments|élémts} $\ne\varnothing$ minimaux de $F_{\mathcal P}$. Les $\mathcal P$ \unsure{correspondent aux} $F$ \uuline{finis}.}

Supposons $X$ un espace topologique. Alors la condition c) équivaut à la suivante : \textcircled{$1^\circ$} $\forall i\in I$, $\overline X_i$ est une réunion des $X_j$ (i.e. est saturée par $R$). \deleted{Montrons} \deleted{\unsure{modération}.} \inserted{La condition d) équivaut à}

\deleted{que cela implique que \textcircled{$2^\circ$} chaque $X_i$ est \uline{loc. fermé} \inserted{dans $X$}, i.e. ouvert dans $\overline X_i$, ou encore que $\overline X_i-X_i$ est fermé. En effet, c'est une réunion des ``feuilles'' $X_j$, avec $j\ne i$, or si (pour \inserted{un} $j\ne i$) ... $X_j\subset\overline X_i$, je dis que l'on a $\overline X_j\cap X_i=\varnothing$ : en effet, je dis que ceci équivaut à :}

\margin{Mais $1^\circ$ + les $X_i$ loc. fermés n'implique \unsure{pas} \deleted{\unsure{ceci}} (\unsure{Ex} : \deleted{\unsure{le}} $X=\{x\}\unsure{\amalg|\cup}(X-\{x\})$, $x$ pt fermé non isolé) \deleted{\unsure{de $X$ et leurs complémentaires $\in F$ (pt. fermé isolé)} $X$ étant séparé}.}

$2^\circ)$ Si $i,j\in I$, alors $\overline X_i=\overline X_j\Longrightarrow X_i=X_j$.

\deleted{\unsure{Néc}.} \uline{Néc} \deleted{\unsure{évident}} \inserted{Comme \unsure{e)|c)} est vérifié trivialement, \unsure{chaque} $X_i$ est \unsure{réunion finie de} parties \unsure{étrangères|étagées} loc. fermées, donc \unsure{l’un au moins est|son intérieur est} $\ne\varnothing$ \unsure{dans} $\overline X_i$.}

\deleted{Si $\overline X_i=\overline X_j\ (=Z)$, comme $X_i$ et $X_j$ sont \unsure{denses} dans $Z$, \unsure{donc} $\unsure{Z}$, \unsure{d'où} $X_i\cap X_j\unsure{\ne}\varnothing$, \unsure{vide}.}

\inserted{\unsure{et} $\unsure{X_j^\circ}$ \unsure{est loc. fermé, donc dense dans} $\overline X_i$. Mais alors $\overline X_i=\overline X_j\Longrightarrow X_i\cap X_j\ne\varnothing$, donc $i=j$, OK.}

Je dis que $1^\circ)$ et $2^\circ)$ impliquent que les $X_i$ sont \uline{loc$^t$ fermés}. En effet,

\mspage{12}
